%% ============================================================
%%  3.  BACKGROUND
%% ============================================================
\section{Background}
\label{sec:background}

\subsection{Deliberative Democracy Axioms}

\citeauthor{diagnostic2026}~(\citeyear{diagnostic2026}) ground multi-agent
debate evaluation in eight axioms drawn from deliberative democracy
theory~\cite{fishkin2009deliberation,arrow1950difficulty,sen1970collective}:

\begin{table}[h]
\centering
\caption{Deliberative axioms and their diagnostic metrics.
This paper treats Non-domination and Pluralism as \emph{enforceable}
constraints, not merely observable outcomes.}
\label{tab:axioms}
\small
\begin{tabular}{lll}
\toprule
\textbf{Axiom} & \textbf{Meaning} & \textbf{Metric} \\
\midrule
Monotonicity        & Meaningful belief change over time      & Engagement \\
Responsiveness      & Updates causally linked to peer arguments & Responsiveness \\
Non-domination      & Influence distributed, not concentrated & Influence Asymmetry \\
Pluralism           & Meaningful disagreement maintained      & Balance / PER \\
Openness            & Engage with counterarguments            & Engagement / PER \\
Cond.\ Convergence  & Consensus from reasoned interaction     & Consensus Reliability \\
Stability           & Outcomes robust to unilateral deviation & Stability \\
Agent Utility       & Decisions do not systematically harm    & Group Welfare \\
\bottomrule
\end{tabular}
\end{table}

Whereas \citeauthor{diagnostic2026} use these metrics for \emph{post-hoc}
evaluation, our key normative move is to treat Non-domination and Pluralism
as procedural inputs that can be actively \emph{enforced} through adaptive
interventions.

\subsection{Diagnostic Metrics}
\label{sec:metrics}

We inherit the six diagnostic metrics from \citeauthor{diagnostic2026} with
their categorical adaptations for multiple-choice answers, and add two new
outcome-level metrics.

\paragraph{Process metrics (existing).}
\emph{Engagement} (Eq.~4 in \citeauthor{diagnostic2026}) measures
quality-weighted stance changes; low engagement signals dogmatism.
\emph{Responsiveness} (Eqs.~5--6) measures quality-weighted movement toward
peer consensus; low responsiveness signals parallel reasoning without
genuine interaction.
\emph{Influence Asymmetry} $\hat{I}$ (Eqs.~6--7) measures normalized
entropy of peer-attribution weights; high $\hat{I}$ signals domination.
\emph{Balance} $B$ (Eqs.~8--11) penalises both premature collapse and
oscillation; low $B$ signals premature convergence.
\emph{Stability} and \emph{Group Welfare} (Eqs.~14--16) measure Nash
equilibrium fraction and average utility gain.

\paragraph{New outcome metrics.}
\emph{Consensus Reliability} (CR) conditions accuracy on consensus:
$\text{CR}=P(\text{correct}\mid\text{consensus reached})$.
Standard accuracy can be gamed by increasing the moderator fallback rate,
which CR controls for.
\emph{Premature Elimination Rate} (PER, defined in \S\ref{sec:formulation})
directly measures the failure mode we target: correct hypotheses that are
abandoned before the final round.
PER is zero only if every correct hypothesis that was ever proposed survives
to the last round.
